% ============================================================
% Capítulo 12 — Métricas: Top-k, AP, mAP, CMC
% ============================================================
\chapter{Métricas: Top-k, AP, mAP y CMC}
\label{cap:metricas}

Este capítulo desglosa línea por línea cada una de las métricas
que usa el proyecto. Si nunca trabajaste con \emph{re-ID} o
recuperación de imágenes, esta es la lectura obligatoria: vas a
entender qué significa cada número del Top-k y cómo se calcula
exactamente.

\section{El problema de evaluación}

El modelo produce, para cada imagen y cada texto, un vector de
1024 dimensiones (un \emph{embedding}). Para evaluar, comparamos
esos vectores con la \textbf{similitud coseno}:

\begin{equation}
  \text{sim}(\mathbf{I}', \mathbf{T}') =
  \frac{\mathbf{I}' \cdot \mathbf{T}'}
       {\|\mathbf{I}'\| \cdot \|\mathbf{T}'\|}
  \label{eq:sim-cos}
\end{equation}

El resultado es una \textbf{matriz de similitud}
$\mathbf{S} \in \mathbb{R}^{N_{\text{img}} \times N_{\text{txt}}}$
donde $S_{ij}$ indica qué tan parecida es la imagen $i$ al
texto $j$.

\subsection*{Galería vs.\ query}

En este proyecto:
\begin{itemize}
  \item \textbf{Query} = descripción textual (texto).
        Son $N_{\text{txt}} = 29\,930$ en el split test.
  \item \textbf{Gallery} = imágenes.
        Son $N_{\text{img}} = 4\,256$ en el split test.
\end{itemize}

Para cada query (texto) queremos encontrar la imagen correcta
en la gallery. El ranking se obtiene ordenando las imágenes por
similitud descendente.

\section{Top-k accuracy (CMC)}

\subsection*{Qué es}

El \emph{Top-k accuracy} (también llamado \textbf{CMC}, Cumulative
Matching Characteristic) responde a la pregunta:

\begin{quote}
\textit{Dada una descripción textual, ¿la imagen correcta está
entre las primeras $k$ del ranking?}
\end{quote}

\noindent\textbf{Ejemplo:} si para el texto ``\emph{mujer con
camiseta negra y pantalones caqui}'' el ranking de imágenes es:
\begin{enumerate}
  \item img\_382 (correcta)
  \item img\_124
  \item img\_991
\end{enumerate}

\noindent entonces Top-1 = 100\% (acertó de entrada), Top-5 = 100\%
(también está en el top-5). Si la imagen correcta hubiera
aparecido en la posición 7, Top-1 = 0\% pero Top-10 = 100\%.

\subsection*{Cómo se calcula}

Para todos los queries a la vez:

\begin{equation}
  \text{Top-}k = \frac{1}{N_q} \sum_{i=1}^{N_q}
                  \mathbb{1}[\text{correcta}_i \leq k]
  \label{eq:topk}
\end{equation}

\noindent donde $\mathbb{1}[\cdot]$ es la función indicadora (1 si
se cumple, 0 si no).

\section{Average Precision (AP)}

\subsection*{Qué es}

El \acronym{AP}{AP} mide la calidad del ranking \emph{completo},
no sólo del top-1. Es el área bajo la curva
\emph{precisión-recall}.

\subsection*{Ejemplo paso a paso}

Supongamos que para una query hay \textbf{3 imágenes correctas}
en la gallery y el modelo las rankea así (de mayor a menor
similitud):

\begin{table}[H]
\centering
\caption{Ejemplo de ranking para una query con 3 positivos.}
\label{tab:ap-ejemplo}
\begin{tabular}{ccccl}
\toprule
\textbf{Pos.} & \textbf{¿Correcta?} & \textbf{Precisión} & \textbf{Recall} & \textbf{Nota} \\
\midrule
1 & \cmark & 1/1 = 1.00 & 1/3 = 0.33 & primer positivo \\
2 & \xmark & 0/2 = 0.00 & 1/3 = 0.33 & falso positivo \\
3 & \cmark & 1/3 = 0.33 & 2/3 = 0.67 & segundo positivo \\
4 & \xmark & 1/4 = 0.25 & 2/3 = 0.67 & falso positivo \\
5 & \cmark & 2/5 = 0.40 & 3/3 = 1.00 & tercer positivo \\
\bottomrule
\end{tabular}
\end{table}

\noindent\textbf{Precisión} = (correctos vistos) / (total vistos).\\[2pt]
\noindent\textbf{Recall} = (correctos vistos) / (total correctos).

El AP es la integral de la curva precisión-recall. Numéricamente,
la regla del trapecio da:

\begin{align}
  \text{AP} &=
    \sum_{i=1}^{N} \Delta\text{recall}_i \cdot
                 \frac{\text{prec}_i + \text{prec}_{i-1}}{2}
\end{align}

\noindent Aplicado a las posiciones donde hay positivos
(1, 3, 5):

\begin{align*}
  \Delta\text{recall}_1 &= 1/3 \quad \text{prec}_1=1, \text{prec}_0=1 \\
  \Delta\text{recall}_3 &= 1/3 \quad \text{prec}_3=1/3, \text{prec}_2=0 \\
  \Delta\text{recall}_5 &= 1/3 \quad \text{prec}_5=2/5, \text{prec}_4=1/4 \\
  \text{AP} &= \frac{1}{3}(1 + \frac{0.5+0}{2}) +
               \frac{1}{3}(\frac{1}{3}+\frac{1}{4}) \cdot \frac{1}{2} \cdot 2 \\
            &= 0.667 + 0.097 + 0.097 \\
            &= 0.86
\end{align*}

\subsection*{Implementación línea por línea}

\begin{minted}[numbers, firstnumber=20, linerange=20-53]{python}
def calculate_ap(similarity, label_query, label_gallery):
    """
    Args.:
        similarity:     vector 1-D con similitud a cada imagen de la gallery
        label_query:    pid (int) del query
        label_gallery:  vector 1-D con pids de la gallery (mismo orden que similarity)
    Returns: (ap, cmc) o (None, None) si no hay positivos
    """

    # Ordenar índices por similitud DESCENDENTE (mejor primero)
    # np.argsort devuelve ascending; [::-1] invierte el orden
    # np es numpy: import numpy as np
    index = np.argsort(similarity)[::-1]

    # good_index: posiciones en la gallery donde el pid == label_query
    # np.argwhere devuelve índices donde se cumple la condición
    good_index = np.argwhere(label_gallery == label_query)

    # cmc: vector de ceros del mismo tamaño que el ranking
    # Lo llenaremos con 1 desde la posición del primer match
    cmc = np.zeros(index.shape)

    # Máscara booleana: True donde index[i] está en good_index
    # np.in1d chequea pertenencia a un conjunto, eficiente
    mask = np.in1d(index, good_index)

    # Posiciones (en el ranking) donde hay match
    # precision_result[i] = posición en el ranking del i-ésimo match
    precision_result = np.argwhere(mask == True)
    precision_result = precision_result.reshape(
        precision_result.shape[0])

    # Si hay al menos un positivo:
    if precision_result.shape[0] != 0:

        # Marcar con 1 desde el primer match en adelante
        # cmc[k] = 1 significa "el match aparece en posición <= k"
        cmc[int(precision_result[0]):] = 1

        # d_recall: cuánto aumenta el recall por cada positivo
        # 1/N donde N = total de positivos en la gallery
        d_recall = 1.0 / len(precision_result)

        # Acumular AP con regla del trapecio
        ap = 0
        for i in range(len(precision_result)):

            # Precisión en la posición del i-ésimo match
            # = (matches hasta aquí) / (rank del match)
            precision = (i + 1) * 1.0 / (precision_result[i] + 1)

            # Precisión del match anterior (para el trapecio)
            if precision_result[i] != 0:
                old_precision = i * 1.0 / precision_result[i]
            else:
                # Caso especial: primer match en posición 0
                old_precision = 1.0

            # Área del trapecio entre (old_prec, precision) con altura d_recall
            ap += d_recall * (old_precision + precision) / 2

        return ap, cmc

    else:
        # No hay positivos para esta query (no debería pasar)
        return None, None
\end{minted}

\subsection*{Detalles a tener en cuenta}

\begin{itemize}
  \item Si la imagen correcta aparece en la posición 1, AP = 1.0
        (perfect ranking).
  \item Si aparece en la última posición, AP es bajo pero no 0
        (porque se premia que aparezca \emph{al final}, no
        totalmente ausente).
  \item Si una query tiene \emph{un solo} positivo en la gallery,
        AP reduce a Top-1 (porque el recall siempre crece igual).
  \item Si tiene varios positivos, AP captura mejor la calidad
        que Top-1.
\end{itemize}

\section{mAP (mean Average Precision)}

\subsection*{Qué es}

\acronym{mAP}{mAP} es simplemente el promedio de AP sobre todos
los queries:

\begin{equation}
  \text{mAP} = \frac{1}{N_q} \sum_{i=1}^{N_q} \text{AP}_i
  \label{eq:map}
\end{equation}

\noindent donde $N_q$ es el número de queries (textos).

\subsection*{Por qué mAP y no sólo Top-1}

\begin{itemize}
  \item \textbf{Top-1} sólo cuenta si el primero es correcto.
        Ignora si los otros positivos están bien rankeados.
  \item \textbf{mAP} considera el ranking \emph{completo}: si una
        query tiene 5 positivos y todos están en los primeros 5,
        AP = 1.0. Si dos están fuera del top-5, AP < 1.0.
\end{itemize}

\noindent\textbf{En este proyecto} cada query (texto) tiene
\emph{un solo} positivo en la gallery (la imagen correspondiente).
En ese caso mAP y Top-1 están correlacionados pero NO son
iguales:
\begin{itemize}
  \item Top-1 cuenta cuántos queries tienen el positivo en
        posición 1.
  \item mAP promedia la calidad del ranking: un query con el
        positivo en posición 2 da AP = 0.5 (en el caso ideal),
        mientras que Top-1 = 0 para ese query.
\end{itemize}

\section{CMC: el vector completo}

\subsection*{Qué es}

El proyecto devuelve \texttt{cmc} que es un vector de tamaño
$N_{\text{gallery}}$. Para cada posición $k$:

\begin{equation}
  \text{cmc}[k] = \frac{1}{N_q}
                  \sum_{i=1}^{N_q} \mathbb{1}[\text{rank}_i \leq k]
  \label{eq:cmc}
\end{equation}

\noindent Es la fracción acumulada de queries cuyo positivo
aparece en posición $\leq k$. \texttt{cmc[0]} = Top-1,
\texttt{cmc[4]} = Top-5, etc.

\subsection*{Implementación}

\begin{minted}[numbers, firstnumber=56, linerange=56-77]{python}
def evaluate(similarity, label_query, label_gallery):
    """
    similarity:     matriz (N_img, N_txt) de similitud coseno
    label_query:    vector (N_txt,) con pids de los textos (queries)
    label_gallery:  vector (N_img,) con pids de las imágenes (gallery)
    Returns: (cmc, mAP)
    """

    # Convertir de torch.Tensor a numpy array
    label_query = label_query.numpy()
    label_gallery = label_gallery.numpy()

    # Vector acumulador de CMC
    # Tamaño = número de imágenes en la gallery
    # Cada posición k almacena cuántos queries acertaron en top-k
    cmc = np.zeros(label_gallery.shape)

    # Acumulador de AP (para luego sacar promedio = mAP)
    ap = 0

    # Para cada query de texto:
    for i in range(len(label_query)):

        # Calcular AP para esta query
        # similarity[:, i] = columna i = similitud de TODAS las imágenes
        #                    con el texto i-ésimo (orden: filas son imágenes)
        ap_i, cmc_i = calculate_ap(
            similarity[:, i],
            label_query[i],
            label_gallery)

        # Acumular cmc (suma de indicadores)
        cmc += cmc_i

        # Acumular AP
        ap += ap_i

    # Normalizar: cmc[k] = fracción de queries que acertaron en top-k
    cmc = cmc / len(label_query)

    # Normalizar: mAP = promedio de AP sobre todos los queries
    map = ap / len(label_query)

    """
    Comentario sobre cmc:
    cmc[i] es un vector [0, 0, ..., 1, 1, ..., 1] donde el primer 1
    está en la posición del primer match correcto. rank-n y todos
    los ranks posteriores suman un acierto. Al sumar todos los
    vectores y dividir por n_query obtenemos el rank-k accuracy.
    """

    return cmc, map
\end{minted}

\section{Interpretación de números}

\subsection*{Resultados típicos para CUHK-PEDES}

\begin{table}[H]
\centering
\caption{Resultados esperados según el README y el paper.}
\label{tab:resultados-tipicos}
\begin{tabular}{lcccc}
\toprule
\textbf{Setup} & \textbf{Top-1} & \textbf{Top-5} & \textbf{Top-10} & \textbf{mAP} \\
\midrule
Sin entrenar (modelo random) & \textasciitilde 0.02\% & 0.05\% & 0.10\% & 0.10\% \\
HF (128×128) 50 batches     & 0.02\% & 0.08\% & 0.17\% & 0.10\% \\
HF entrenado completo       & 45--52\% & \textasciitilde 70\% & \textasciit\sim 80\% & \textasciitilde 40\% \\
Original paper (paper)      & 54.02\% & 75.51\% & 85.28\% & \textasciitilde 50\% \\
\bottomrule
\end{tabular}
\end{table}

\subsection*{Cómo interpretar un resultado ``Top-1 = 50\%''}

Dice que, dado un texto al azar del test set, en el 50\% de los
casos la imagen correcta aparece en la primera posición del
ranking.

\begin{itemize}
  \item Es ``bueno'' para ReID por texto (un humano也很难
        hacerlo mejor).
  \item Es ``malo'' comparado con face recognition ($>$99\%).
  \item El gap entre Top-1 y Top-5 indica cuántas ``casi
        correctas'' hay en el top-5: si Top-5 = 75\% y
        Top-1 = 50\%, el 25\% restante tiene el correcto en
        posiciones 2--5.
\end{itemize}

\subsection*{Cómo interpretar mAP}

\begin{itemize}
  \item mAP es siempre \textbf{menor o igual} a Top-1
        (matemáticamente: $\text{mAP} \leq \text{cmc}[0]$).
  \item Si Top-1 $\approx$ mAP, el modelo es ``perfecto'' en
        sus predicciones (los positivos casi siempre están en
        posición 1).
  \item Si Top-1 $>$ mAP $\cdot 2$, hay muchas predicciones donde
        el correcto está en posiciones 2--5 (modelo ``casi
        bueno'').
\end{itemize}

\section{Ejemplo numérico completo}

Supongamos 3 queries y 5 imágenes en gallery. La matriz de
similitud es:

\begin{verbatim}
       Q1    Q2    Q3
img1 [ 0.9   0.1   0.3 ]
img2 [ 0.5   0.8   0.2 ]
img3 [ 0.2   0.6   0.7 ]
img4 [ 0.7   0.3   0.5 ]
img5 [ 0.4   0.5   0.4 ]
\end{verbatim}

\noindent y los pids (queries vs gallery):
\begin{itemize}
  \item Q1 busca img2 (pid=2)
  \item Q2 busca img3 (pid=3)
  \item Q3 busca img3 (pid=3)
\end{itemize}

\subsection*{Calcular Top-1}

Para cada query ordenamos la columna de similitud:

\begin{itemize}
  \item \textbf{Q1}: img1 (0.9), img4 (0.7), img2 (0.5), img5 (0.4), img3 (0.2).
        Correcto (img2) en posición \textbf{3}. Top-1: NO.
  \item \textbf{Q2}: img2 (0.8), img3 (0.6), img5 (0.5), img4 (0.3), img1 (0.1).
        Correcto (img3) en posición \textbf{2}. Top-1: NO.
  \item \textbf{Q3}: img3 (0.7), img4 (0.5), img1 (0.3), img2 (0.2), img5 (0.4)
        Correcto (img3) en posición \textbf{1}. Top-1: SÍ.
\end{itemize}

\noindent\textbf{Top-1 = 1/3 = 33.3\%}.

\subsection*{Calcular mAP}

\begin{itemize}
  \item \textbf{AP(Q1)}: positivo en posición 3 de 5.
        AP = 1/3 $\cdot$ (1 + (1/3+0)/2) = 1/3 + 1/9 = 4/9 $\approx$ 0.44.
  \item \textbf{AP(Q2)}: positivo en posición 2 de 5.
        AP = 1/1 $\cdot$ (1 + 0)/2 = 0.5.
  \item \textbf{AP(Q3)}: positivo en posición 1 de 5.
        AP = 1.
\end{itemize}

\noindent\textbf{mAP = (0.44 + 0.5 + 1) / 3 $\approx$ 0.65}.

\subsection*{Calcular Top-5}

Todos los positivos están en posiciones $\leq$ 5 (trivialmente),
así que \textbf{Top-5 = 100\%}.

\section{Resumen}

\begin{table}[H]
\centering
\caption{Resumen de las cuatro métricas.}
\label{tab:resumen-metricas}
\begin{tabular}{ll}
\toprule
\textbf{Métrica} & \textbf{Qué captura} \\
\midrule
Top-1 (cmc[0]) & ¿La imagen correcta es la \#1? \\
Top-5 (cmc[4]) & ¿Está en el top 5? \\
Top-10 (cmc[9]) & ¿Está en el top 10? \\
mAP & Calidad global del ranking \\
\bottomrule
\end{tabular}
\end{table}

\noindent\textbf{Consejo para reportar resultados:} siempre
reportar al menos Top-1, Top-5, Top-10 \emph{y} mAP. Top-1 solo
es insuficiente porque no captura la calidad del ranking más
allá del primer lugar.
